\documentclass[11pt,a4paper]{article}
\usepackage[utf8]{inputenc}
\usepackage[ukrainian]{babel}
\usepackage[T1]{fontenc}
\usepackage{hyperref}
\usepackage{enumitem}
\usepackage{longtable}
\usepackage{fancyhdr}
\usepackage{sectsty}

\title{ТЕХНІЧНЕ ЗАВДАННЯ \\ на створення \\ комплексної системи захисту інформації \\ в автоматизованій системі ERP/1}
\author{На основі НД ТЗІ 3.7-001-99}
\date{\today}

\begin{document}

\maketitle

\section{Загальні відомості}
Повне найменування КСЗІ: Комплексна система захисту інформації в автоматизованій системі ERP/1 (КСЗІ ERP/1). 

Умовне позначення: КСЗІ-ERP/1.

Шифр теми: [Вказати шифр].

Найменування підприємств-розробників та замовника: [Вказати реквізити].

Перелік документів, на підставі яких створюється КСЗІ: 
\begin{itemize}
    \item НД ТЗІ 3.7-001-99 «Методичні вказівки щодо розробки технічного завдання на створення комплексної системи захисту інформації в автоматизованій системі»;
    \item НД ТЗІ 2.5-005-99 «Класифікація автоматизованих систем...»;
    \item НД ТЗІ 2.5-004-99 «Критерії оцінки захищеності...»;
    \item Інші нормативні документи.
\end{itemize}

Планові терміни початку і закінчення робіт: [Вказати].

Джерела фінансування: [Вказати].

Порядок оформлення результатів: [Вказати].

\section{Мета і призначення комплексної системи захисту інформації}
Метою розробки КСЗІ в АС ERP/1 є забезпечення захисту інформації з обмеженим доступом від несанкціонованого доступу та витоку технічними каналами відповідно до вимог законодавства України та нормативних документів Департаменту спеціальних телекомунікаційних систем та захисту інформації СБ України.

Функціональне призначення: Реалізація комплексного захисту в ERP-системі, яка обробляє критичну бізнес-інформацію, фінансові дані, персональні дані тощо.

\section{Загальна характеристика автоматизованої системи та умов її функціонування}
АС ERP/1 призначена для автоматизації бізнес-процесів підприємства (фінанси, логістика, HR тощо).

Структурна схема: Клієнт-серверна архітектура з центральним сервером баз даних, веб-клієнтами, інтеграціями з зовнішніми системами.

Клас АС: Відповідно до НД ТЗІ 2.5-005-99 — [Вказати клас, наприклад 2Б або 1Г].

Характеристики інформації: Обробляється інформація з обмеженим доступом (конфіденційна, комерційна таємниця).

Персонал: [Кількість користувачів, ролі, рівні доступу].

Фізичне середовище: Розміщення в [описати].

Потенційні загрози: Несанкціонований доступ, витік через ПЕМВН, внутрішні загрози тощо.

Функціонуючі засоби захисту: [Описати наявні].

\section{Вимоги до комплексної системи захисту інформації}

\subsection{Вимоги в частині захисту від несанкціонованого доступу}
Функціональний профіль захищеності: [Вказати профіль відповідно до НД ТЗІ 2.5-005-99].

Політика безпеки:
\begin{itemize}
    \item Об'єкти доступу: Бази даних, файли, модулі ERP.
    \item Принципи керування доступом: Рольовий доступ (RBAC), мандатний контроль.
    \item Правила розмежування потоків.
    \item Аудит та реєстрація подій.
\end{itemize}

Послуги безпеки (відповідно до НД ТЗІ 2.5-004-99):
\begin{enumerate}
    \item Конфіденційність — рівень [вказати].
    \item Цілісність — рівень [вказати].
    \item Доступність — рівень [вказати].
    \item Спостереженість — рівень [вказати].
\end{enumerate}

Рівень гарантій: [Вказати, наприклад EAL3 або відповідний].

\subsection{Вимоги в частині захисту від витоку інформації технічними каналами}
Загальні вимоги: Захист від ПЕМВН відповідно до ТР ЕОТ-95 та ТР ПЕМВН-95.

Показники захищеності:
\begin{itemize}
    \item Нормовані значення співвідношення сигнал/шум.
    \item Застосування фільтрів, зашумлення, екранування.
\end{itemize}

Спецперевірки та спецдослідження ЗОТ.

\section{Вимоги до складу проектної та експлуатаційної документації}
\begin{itemize}
    \item Технічне завдання на КСЗІ.
    \item Проектна документація (архітектура, специфікації).
    \item Експлуатаційна документація (інструкції, керівництва).
    \item Програми та методики випробувань.
    \item Акти приймання.
\end{itemize}

\section{Етапи виконання робіт}
\begin{enumerate}
    \item Попередній етап (аналіз, класифікація).
    \item Проектування та розробка КСЗІ.
    \item Випробування та введення в експлуатацію.
\end{enumerate}

Календарний план: [Вказати].

\section{Порядок внесення змін і доповнень}
Зміни оформлюються доповненням, погоджуються та затверджуються в установленому порядку.

\section{Порядок проведення випробувань}
Розробка Програми та методики випробувань. Проведення комісією. Використання умовної інформації. Документи завершення: акт, сертифікат, наказ.

\end{document}
